% Generated by Hypo-LaTeX. Do not edit directly.
\documentclass{book}
\usepackage{hypolatex}
\HypoSetAnswerMode{student}

\HypoSetMetadata{title}{Review Question Variants}
\HypoSetMetadata{author}{Test Worker}
\HypoSetMetadata{theme}{plain}
\HypoSetMetadata{layout}{standard}

\begin{document}
\HypoDocumentStart

\chapter{Review Sheet}

Use the same source for student, review, and teacher copies.

\section{Card Question}

\begin{HypoQuestion}[rq:power-rule][题目：Power Rule][][card]
What is the derivative of \(x^2\)?

\begin{itemize}
\tightlist
\item
Include the coefficient and simplify the result.
\end{itemize}
\end{HypoQuestion}

\begin{HypoHint}
Hint text: Use the power rule before simplifying.
\end{HypoHint}

\begin{HypoAnswer}
Official answer text: The derivative is two x.
\end{HypoAnswer}

\begin{HypoSolution}
Official solution text: Bring down the exponent and subtract one.
\end{HypoSolution}

\section{Plain Question}

\begin{HypoQuestion}[rq:plain-discussion][非 Card 开放推理题][][plain]
Explain why the derivative rule is useful before computing a final numeric answer.

\begin{itemize}
\tightlist
\item
State the rule in one sentence.
\item
Mention one common mistake.
\end{itemize}
\end{HypoQuestion}

\begin{HypoHint}[提示][plain]
Use the rule name before doing algebra.
\end{HypoHint}

\begin{HypoAnswer}[参考答案][plain]
It turns exponents into quick calculations.
\end{HypoAnswer}

\begin{HypoSolution}[解析][plain]
Labels separate the semantic parts without boxes.
\end{HypoSolution}
\end{document}
